\documentclass[11pt]{article}

\usepackage[margin=1in]{geometry}
\usepackage{amsmath,amssymb,amsthm,bm}
\usepackage{enumitem}
\usepackage{hyperref}
\usepackage{xcolor}

\hypersetup{
  colorlinks=true,
  linkcolor=blue!55!black,
  citecolor=blue!55!black,
  urlcolor=blue!55!black
}

\newcommand{\R}{\mathbb{R}}
\newcommand{\E}{\mathbb{E}}
\newcommand{\Prob}{\mathbb{P}}
\newcommand{\Normal}{\mathcal{N}}
\newcommand{\Orth}{\mathcal{O}}
\newcommand{\argmax}{\operatorname*{arg\,max}}
\newcommand{\argmin}{\operatorname*{arg\,min}}
\newcommand{\diag}{\operatorname{diag}}
\newcommand{\rank}{\operatorname{rank}}
\newcommand{\PhiN}{\Phi}
\newcommand{\independent}{\mathrel{\perp\!\!\!\perp}}

\title{A Likelihood-Based Pretraining and Refinement Algorithm\\
for Probit Independent-Mixture Factor Models}
\author{}
\date{\today}

\begin{document}
\maketitle

\section{Why Estimate the Signal Before Augmenting?}

The key advantage of the likelihood-based pretraining step is that it targets
the identifiable signal first.  In the probit factor model,
\[
  \Prob(X_{ij}=1)=\PhiN(\alpha_j+B_{ij}),
  \qquad
  B=F\Lambda^\top,
\]
the systematic low-rank object is \(B\).  By contrast, an augmented latent
Gaussian draw satisfies
\[
  Z_{ij}=\alpha_j+B_{ij}+\epsilon_{ij}.
\]
Thus an SVD of an augmented \(Z\) matrix is trying to recover the low-rank
signal from a noisy latent realization.  Even if the augmentation is correctly
specified, the sampled noise can perturb the singular vectors, introduce
Monte Carlo variation, and slow the initialization.

The alternative is to estimate \(B\) directly from the binary probit
likelihood and then take the SVD of \(\widehat B\).  This gives a cleaner
initial estimate of the signal subspace, avoids repeated stochastic
augmentation during pretraining, and provides a more direct theoretical route:
if \(\widehat B\) consistently estimates the low-rank response surface, then
its singular subspace consistently estimates the latent factor subspace.  The
subsequent rotation and refinement stages then operate on a likelihood-based
estimate of the systematic signal rather than on noisy augmented draws.

\section{Goal}

The main goal of pretraining is not to estimate uniquely identified factors
and loadings.  It is to estimate the low-rank probit response surface.  We
write
\[
  \Theta = \bm 1_n\alpha^\top + B,
  \qquad
  B=F\Lambda^\top,\qquad \rank(B)\leq H.
\]
The product \(B\) is identified by the binary likelihood, while a particular
factorization \(B=F\Lambda^\top\) is not unique.  This distinction is the
reason for the staged strategy: first estimate \(B\) and its subspace, then
rotate that subspace into mixture-factor coordinates, and finally refine the
full structured model.

\section{Stage 1: Estimate the Low-Rank Signal}

For a supplied or selected working rank \(H\), the first stage estimates
\((\alpha,B)\) by maximizing the rank-constrained probit likelihood
\begin{equation}
  (\widehat\alpha,\widehat B)
  \in
  \argmax_{\alpha,\,B:\ \rank(B)\leq H,\; \bm 1_n^\top B=0}
  \ell_H(\alpha,B),
  \label{eq:B-objective}
\end{equation}
where
\begin{equation}
  \ell_H(\alpha,B)=
  \sum_{i=1}^n\sum_{j=1}^p
  \left[
    X_{ij}\log\PhiN(\alpha_j+B_{ij})
    +(1-X_{ij})\log\{1-\PhiN(\alpha_j+B_{ij})\}
  \right].
\end{equation}
The centering constraint separates item intercepts from the low-rank
interaction.  Under mild regularity conditions, this likelihood-based step
recovers the population rank-\(H\) signal subspace.  Taking the SVD of
\(\widehat B\) then gives a consistent initialization of the left singular
subspace; when the working rank matches the data-generating rank, this is the
true factor-score subspace.

Computationally, this can be done with a simple deterministic EM-SVD
algorithm.  Given the current low-rank signal, the E-step computes
\[
  W_{ij}=\E\{Z_{ij}\mid X_{ij},\alpha_j+B_{ij}\},
\]
using the usual truncated-normal probit formula.  The M-step projects \(W\)
back onto the constrained low-rank class: set
\(\widehat\alpha=\operatorname{colMeans}(W)\), center \(W\) by columns, and
take the rank-\(H\) SVD.  Thus the rank constraint is enforced by SVD
truncation and the intercept constraint is enforced by column-centering.

This removes the need to use stochastic augmented-\(Z\) draws as the
pretraining estimator.  The latent-Gaussian representation is still useful for
EM calculations, but the target is the observed binary likelihood in
\eqref{eq:B-objective}.  In practice this is faster because pretraining no
longer repeatedly samples full \(n\times p\) augmented matrices and recomputes
the rotation from noisy draws.  Instead, each EM iteration uses deterministic
conditional means followed by a low-rank SVD projection.

\section{Stage 2: Rotate the Subspace}

After estimating \(\widehat B\), we take a convenient factorization, for
example from the SVD,
\[
  \widehat B=\widehat F\widehat\Lambda^\top.
\]
These scores estimate the correct subspace, but their axes are arbitrary up
to rotation.  The mixture assumption supplies the missing orientation: we look
for the rotation under which the score coordinates look like independent
one-dimensional mixtures.  This is the ICA-like step of the algorithm.

\section{Stage 3: Refine the Structured Model}

Starting from the rotated scores, the refinement stage updates the full
product-mixture factor model:
\[
  \alpha,\quad F,\quad \Lambda,\quad
  \{\pi_{hg},\mu_{hg},\sigma^2_{hg}\}.
\]
In our current implementation, refinement uses MAP mixture updates and can
include sparsity penalties on the loadings.  The practical point is that
refinement starts from a likelihood-based estimate of the rank-\(H\) subspace,
rather than from a noisy sampled augmentation.  The refinement step is where
the model moves from a good low-rank probit representation to the full
independent-mixture factor model used for interpretation and prediction.

\paragraph{Reporting convention.}
The final fitted model is still invariant to location, scale, sign, and column
ordering conventions.  After convergence, we apply a canonical normalization
only for reporting.  If \(m_h\) and \(s_h\) are the fitted marginal mixture
mean and standard deviation for factor \(h\), set
\[
  f_{ih}^\star=\frac{f_{ih}-m_h}{s_h},
  \qquad
  \lambda_{jh}^\star=s_h\lambda_{jh},
  \qquad
  \alpha_j^\star=\alpha_j+\sum_h\lambda_{jh}m_h.
\]
This preserves the probit linear predictor:
\[
  \alpha_j+\bm f_i^\top\bm\lambda_j
  =
  \alpha_j^\star+(\bm f_i^\star)^\top\bm\lambda_j^\star.
\]
So the fitted probabilities do not change; only the parameter convention
changes.

\section{Code Pointers}

The main implementation is \path{R/probit_ifa_em_svd_pretraining.R}, which
contains the low-rank EM-SVD fit, the mixture-rotation pretrainer, and the
pretrain-then-refine wrapper.  The simulation checks are in
\path{scripts/sample_size/}: \path{test_em_svd_pretraining_simulation.R} and
\path{test_final_canonical_normalization.R}.

\end{document}
